\documentclass[11pt,a4paper]{article}
\usepackage[utf8]{inputenc}
\usepackage[T1]{fontenc}
\usepackage{lmodern,amsmath,amssymb,amsthm,mathtools,booktabs,enumitem,microtype}
\usepackage[margin=1in]{geometry}
\usepackage[round,authoryear]{natbib}
\usepackage{xr-hyper}
\usepackage{hyperref}
\externaldocument[SM-]{02_seldon_manifolds_theory_referenced}
\externaldocument[GE-]{01_law_of_graphic_equalisation_theory_referenced}
\externaldocument[NM-]{03_law_of_narrative_momentum_theory_referenced}
\externaldocument[NV-]{04_narrativium_theory_referenced}
\externaldocument[RI-]{11_relational_information_ethics}

\newtheorem{definition}{Definition}[section]
\newtheorem{proposition}[definition]{Proposition}
\newtheorem{corollary}[definition]{Corollary}
\newtheorem{hypothesis}[definition]{Hypothesis}
\theoremstyle{remark}
\newtheorem{remark}[definition]{Remark}

\newcommand{\Frame}{\mathcal F}
\newcommand{\Seldon}{\mathcal S}
\newcommand{\Target}{\mathcal T}
\newcommand{\Eligible}{\mathcal E}
\newcommand{\Pareto}{\mathcal P}
\newcommand{\Meta}{\operatorname{MetaAnalyse}}
\newcommand{\Cohere}{\operatorname{Cohere}}
\newcommand{\Surprise}{\operatorname{Surprise}}

\title{Seldon Manifolds --- Decision Making\\\large Operative Continuation, Prospective Deliberation, and Manifold Revision}
\author{Matthew Parslow\\Independent Researcher}
\date{\textbf{Draft \texorpdfstring{\textemdash}{---}\ \PaperDraft}\\[2pt]\PaperDisplayDate \texorpdfstring{\textbullet}{\textperiodcentered}\ document revision \PaperRevision}

\usepackage{fancyhdr}
\pagestyle{fancy}
\fancyhf{}
\fancyfoot[C]{\thepage}
\fancyfoot[R]{\scriptsize\textsc{\PaperDraft\,r\PaperRevision}}
\renewcommand{\headrulewidth}{0pt}
\fancypagestyle{plain}{\fancyhf{}\fancyfoot[C]{\thepage}%
  \fancyfoot[R]{\scriptsize\textsc{\PaperDraft\,r\PaperRevision}}\renewcommand{\headrulewidth}{0pt}}
% GENERATED by tools/build_paper_metadata.py from paper_metadata.json.
% DO NOT EDIT: edit paper_metadata.json, then regenerate this file.
\providecommand{\PaperDraft}{draft-2026-09-06.1}
\providecommand{\PaperDisplayDate}{6 September 2026}
\expandafter\def\csname PaperRevision@00_is\endcsname{129}
\expandafter\def\csname PaperRevision@01_law_of_graphic_equalisation_theory_referenced\endcsname{294}
\expandafter\def\csname PaperRevision@01a_graphic_equalisation_physics_grounding\endcsname{47}
\expandafter\def\csname PaperRevision@01b_graphic_equalisation_mind_and_self_grounding\endcsname{47}
\expandafter\def\csname PaperRevision@01c_graphic_equalisation_cosmology_grounding\endcsname{11}
\expandafter\def\csname PaperRevision@01d_graphic_equalisation_born_recovery\endcsname{6}
\expandafter\def\csname PaperRevision@01e_graphic_equalisation_communication_grounding\endcsname{6}
\expandafter\def\csname PaperRevision@01f_graphic_equalisation_qualia_and_consciousness_grounding\endcsname{3}
\expandafter\def\csname PaperRevision@01g_graphic_equalisation_total_perspective_vortex_grounding\endcsname{4}
\expandafter\def\csname PaperRevision@01h_graphic_equalisation_information_gravity_quantum_reconciliation\endcsname{5}
\expandafter\def\csname PaperRevision@01i_graphic_equalisation_passive_continuation_momentum_gravity\endcsname{12}
\expandafter\def\csname PaperRevision@01s_graphic_equalisation_statement\endcsname{294}
\expandafter\def\csname PaperRevision@02_seldon_manifolds_theory_referenced\endcsname{31}
\expandafter\def\csname PaperRevision@02a_seldon_manifolds_psychological_grounding\endcsname{15}
\expandafter\def\csname PaperRevision@02b_seldon_manifolds_social_and_psychological_application\endcsname{7}
\expandafter\def\csname PaperRevision@02c_seldon_manifolds_decision_making\endcsname{11}
\expandafter\def\csname PaperRevision@03_law_of_narrative_momentum_theory_referenced\endcsname{32}
\expandafter\def\csname PaperRevision@04_narrativium_theory_referenced\endcsname{23}
\expandafter\def\csname PaperRevision@05_frames_relation_and_the_empty_limit\endcsname{11}
\expandafter\def\csname PaperRevision@06_assessing_an_account_of_this_kind\endcsname{29}
\expandafter\def\csname PaperRevision@07_causal_attitude_networks_on_the_substrate\endcsname{4}
\expandafter\def\csname PaperRevision@08_certified_answers_as_an_audit\endcsname{4}
\expandafter\def\csname PaperRevision@09_adaptability\endcsname{94}
\expandafter\def\csname PaperRevision@09a_adaptability_implementation\endcsname{7}
\expandafter\def\csname PaperRevision@10_applied_graphic_equalisation\endcsname{7}
\expandafter\def\csname PaperRevision@10a_red_queen_race_civilisational_rate_of_change_llms_biological_timescales_anthropocene_transition\endcsname{13}
\expandafter\def\csname PaperRevision@11_relational_information_ethics\endcsname{16}
\expandafter\def\csname PaperRevision@A1_constructive_graphic_equalisation_implementation_referenced\endcsname{8}
\expandafter\def\csname PaperGenerated@01s_graphic_equalisation_statement\endcsname{1}
\edef\PaperMetadataStem{\jobname}
\expandafter\ifx\csname PaperRevision@\PaperMetadataStem\endcsname\relax
  \PackageError{paper_metadata}{No metadata for \PaperMetadataStem}{Add the paper stem to paper_metadata.json.}
  \def\PaperRevision{?}
\else
  \expandafter\let\expandafter\PaperRevision\csname PaperRevision@\PaperMetadataStem\endcsname
\fi
\newcommand{\PaperDateLine}{\textbf{Draft \texorpdfstring{\textemdash}{---}\ \PaperDraft}\\[2pt]\PaperDisplayDate\ \texorpdfstring{\textbullet}{\textperiodcentered}\ document revision \PaperRevision}
\newcommand{\PaperStatementDateLine}{\textbf{Draft \texorpdfstring{\textemdash}{---}\ \PaperDraft}\\[2pt]\PaperDisplayDate\ \texorpdfstring{\textbullet}{\textperiodcentered}\ statement}
\expandafter\ifx\csname PaperGenerated@\PaperMetadataStem\endcsname\relax
  \newcommand{\PaperFooterStamp}{\textsc{\PaperDraft\,r\PaperRevision}}
\else
  \newcommand{\PaperFooterStamp}{\textsc{\PaperDraft\ / statement / from long r\PaperRevision}}
\fi
\newcommand{\PaperNoticeStamp}{\textbf{DRAFT --- \PaperDraft\ r\PaperRevision} (\PaperDisplayDate).}


\begin{document}
\maketitle
\begin{center}
\fbox{\parbox{0.92\textwidth}{\small
\textbf{DRAFT --- \PaperDraft r\PaperRevision} (\PaperDisplayDate).
Circulated for comment, not as a finished result. This is the decision-making
companion to \emph{Seldon Manifolds --- Theory} (referred to below as
\emph{Seldon}) and inherits its definitions, together with those of
\emph{The Law of Graphic Equalisation --- Theory}, the momentum paper, and
\emph{Narrativium}. This paper proposes an ontological interpretation of
realised decision as embedded continuation, a finite specification of
minimum-cost dynamics, and an empirical hypothesis about deliberation as
restructuring of operative future geometry. These have distinct burdens of
support. The finite construction supplies an implementable instance; it does
not establish that human, organisational, or artificial agents universally
instantiate its cost rule. Operational commitments are intended points of
attack.}}
\end{center}
\vspace{1em}

\begin{abstract}
At a declared agent and action scale, a realised decision is the operative
system's action-bearing continuation through its reachable future geometry.
Deliberation is a present process concerning unrealised continuations that can
change the conditions of subsequent action. This identity places deciding
within the system's organisation. The narrower deliberative use of
``decision'' attributes an action or commitment to cognitive manifold
deformation; that causal attribution is inferred, rather than another event
added to the realised continuation. It applies the series' principle that
provenance is inferred from present traces rather than inherently labelled
in an output. A bare set of reachable histories supplies
no successor rule; an operative description must also include admissibility,
cost fields, standing structure, capacity, and dynamics.

The proposed minimum-continuation hypothesis identifies action with the
lowest effective-cost continuation in the operative comparison domain. A
finite specification makes that hypothesis implementable while separating
local edge cost, state potential, accumulated expenditure, and Narrativium.
Local minimisation does not establish a globally cheapest future. Costs and
dynamics must be grounded independently of observed choices.

Biological evolution supplies the central structural example: variation,
inheritance, and differential reproduction change the organisation that
generates subsequent continuations without requiring prospective modelling.
Autonomous processing is the installed response organisation operating;
hypothetical processing is embedded modelling that can restructure it.
This connects to Type 1/Type 2 cognition without equating all continuation
with Type 1 or all revision with Type 2. The represented decision grammar
retains Narrative, surprise, paid meta-analysis, acceptable target shapes,
and present-state authorisation of self-change. Worked finite examples
separate representation from operative revision and actual consequences.
The account is agent-neutral and proposes independently testable groundings;
it claims neither empirical validation nor a universal energy-minimisation
theorem.
\end{abstract}

\section{Scope: decision as operative continuation}\label{sec:scope}

\emph{Seldon} defines \(\Seldon_{h,\rho}(\Frame)\), the structured representation
at horizon \(h\) and resolution \(\rho\) of future histories reachable under
represented admissible interventions and boundary conditions
(Definition~\ref{SM-def:seldon}). It distinguishes a local admissibility
condition from non-zero-horizon future geometry
(Proposition~\ref{SM-prop:local-vs-manifold}), and represents goals, where
held, as acceptable subsets (Definition~\ref{SM-def:acceptable}). None of
those constructions selects an action. This paper adds an account of the
operative organisation whose continuation realises action, and of present
processes that change its subsequent future geometry.

\begin{definition}[Realised decision]\label{def:realised-decision}
For a declared agent boundary, action scale, and causal resolution, a realised
decision is the action-bearing transition or history of that embedded system.
A represented selection or commitment is a present state that can condition
such a transition; its intended action may remain unrealised. The eventual
outcome is the consequence of the realised history under its boundary
conditions. These uses of decision, commitment, and outcome are kept distinct.
\end{definition}

\begin{remark}[Identity and scope]\label{rem:identity}
The proposed identity is that the realised decision \emph{is} this
continuation. The action-bearing organisation, including any selection
machinery, belongs to the system. The definition does not imply that every
physical transition is a decision; the instance must declare an agent and
action scale. It requires neither conscious endorsement nor an observer to
make the transition occur. The identity alone does not predict which action
occurs, or establish a minimum-cost law.
\end{remark}

\begin{definition}[Deliberative decision attribution]\label{def:decision-attribution}
In the narrower deliberative sense, calling an action or commitment a
decision attributes causal influence to a deliberate cognitive process that
deformed operative future geometry. The process can change a represented
landscape, comparison field, constraint, or installed response organisation.
The attribution is a model of that causal history, inferred from available
evidence. It is distinct from the action-bearing continuation and from the
cognitive process, both of which occur whether or not the attribution is made.
\end{definition}

This distinguishes the paper's broad structural use of realised decision from
the familiar deliberative use of the word. In the latter, the word identifies
how a continuation was influenced. It does not introduce a separate act of
choosing after the cognitive deformation. A decision about a future action
can also name the resulting present commitment; the intended action can
remain unrealised.

The same visible action can arise from an installed autonomous response or
from a cognitively revised organisation. Action alone therefore does not
identify deliberative provenance. Behavioural history, process measurements,
and self-report can constrain the attribution; a self-model is not guaranteed
complete access to its own causes. Evidence of prior thought alone also does
not establish that it influenced the action. The causal question concerns
operative influence, without requiring that removing the deliberation must
always change the outcome: other sufficient causes can yield the same action.

\begin{remark}[Decision provenance is inferred, not inherent]\label{rem:decision-provenance}
The deliberative attribution is an application of the series' provenance
principle. A response carries no self-authenticating certificate, independent
of its supporting relations, identifying it as deliberate, habitual, evolved,
trained, or mechanically conditioned. Explicit annotations or authenticated
records can nevertheless constrain local attribution under declared recording
and trust assumptions. Those provenance terms describe causal relations
inferred from the structure and evidence
available to a modelling frame. The actual causal relations do not depend on
that inference; the inference can be accurate, incomplete, or mistaken.
This inherits Graphic Equalisation's foundational definition of relational
provenance (Definition~\ref{GE-def:provenance}) and its evidence-relative
recovery proposition (Proposition~\ref{GE-prop:provenance-inference}). The
operator-supply result of Proposition~\ref{GE-prop:leak} is another application,
rather than the source of this general principle.

The underlying idea is established source-monitoring and identification
work, not a new discovery by this paper. In human source monitoring, origins
are attributed through memory characteristics and judgment processes, which
can operate automatically or deliberately \citep{johnson1993}. Graphic
Equalisation provides a derivation conditional on its relational commitments
and the declared evidence model;
the present paper applies it to the causal attribution of deliberation.

Present traces and records can strongly constrain a local causal attribution.
Their evidential value does not turn provenance into an additional property
that an output exposes independently of a model. A verbal explanation is
also present structure whose relation to an earlier action must be inferred.
The same response can arise through different processes, and the agent's
self-model is one modelling frame among others. Self-report therefore does
not close the provenance question merely by coming from the actor.
\end{remark}

The argument therefore has three levels. The ontological interpretation
identifies realised decision with embedded continuation. The formal
specification supplies explicit state, fields, and dynamics for a finite
instance. The empirical hypothesis asks whether independently grounded
instances explain persistence and revision in actual agents. A correct finite
construction establishes the second level under its premises; human
application needs additional evidence. The paper supplies no neural mechanism
or complete account of motor execution. Inference, carrying a policy, and
realising an outcome remain different processes within the same system.

\section{The decision frame}\label{sec:frame}

\begin{definition}[Decision frame]\label{def:decision-frame}
For a current frame \(\Frame\), let its decision description be
\[
D(\Frame)=(A_{\Frame},L_{\Frame},S_{\Frame},\Gamma_{\Frame},N_{\Frame},V_{\Frame},R_{\Frame}),
\]
where \(A_{\Frame}\) is aperture, \(L_{\Frame}\) the represented landscape,
\(S_{\Frame}\) the standing Narrative, \(\Gamma_{\Frame}\) represented
admissibility, \(N_{\Frame}\) Narrativium, \(V_{\Frame}\) active comparison
descriptions, and \(R_{\Frame}\) represented residual.
\end{definition}

Aperture and residual are inherited from Graphic Equalisation
(Definitions~\ref{GE-def:aperture} and \ref{GE-def:residual}). In particular,
an aperture omission is absent to the frame, rather than automatically a
residual. Narrativium is inherited as frame-relative capacity for consequential
deformation (Definition~\ref{NV-def:narrativium}); it is a state measurement
and feasibility relation, not expenditure, energy, a conserved balance, or a
universal utility.

\begin{definition}[Represented decision domain]\label{def:domain}
The domain of represented decision evaluation is \(L_{\Frame}\) only. A candidate is not
silently assessed as a negative option when it is outside \(A_{\Frame}\) or
absent from \(L_{\Frame}\).
\end{definition}
This restriction concerns in-frame evaluation. Omitted consequences can still
occur, and unrepresented influences can still change actual continuation.

\begin{definition}[Operative description]\label{def:operative}
At a declared boundary and resolution, an operative description contains state
\(X\), a counterfactual transition relation \(\mathcal R\), admissibility,
standing structure, capacity, comparison fields, and a continuation rule
\(b\). Its enriched future geometry is the reachable family together with
these operative conditions. A bare reachable set is one projection of it.
The agent's represented description \(D(\Frame)\), an observer's measurement
of the operative system, and the actual system are distinct.
\end{definition}

A supplied deterministic grounding can be written
\[
X_{t+1}=T(X_t,B_t),\qquad
D(\Frame_t)=\Pi_{\Frame_t}(X_t,B_t),
\]
where \(X_t,B_t\) are sufficient for that instance and \(B_t\) states relevant
containing conditions. The index orders local transitions; it introduces no
global clock. The projection is not an additional chooser. If the measured
state is insufficient, the observer must use a declared transition relation
or an epistemic probability model, rather than report a known unique
successor. Such uncertainty does not add stochasticity to the inherited
Graph Law.

\begin{remark}[Complete fixing and nontrivial comparison]\label{rem:singleton}
When causally complete antecedents, interventions, and containing conditions
are fixed, the inherited deterministic construction has one continuation
(\emph{Seldon}, Corollary~\ref{SM-cor:singleton}). A minimum over that singleton
has no explanatory force. Nontrivial comparisons concern the family indexed
by counterfactual interventions or a partial description that does not fix
all determining organisation. Actual mechanisms generating the intervention
belong to the operative system. A cost rule that explains their continuation
is an additional hypothesis about that family, not a consequence of
determinism.
\end{remark}

Changing a comparison field while keeping the available actions fixed changes
the enriched operative geometry, but not the bare reachable set. Changing
the set through a constraint is a different deformation. The claim that
changing decision changes the manifold uses this enriched meaning. An
observer's redescription without causal interaction need not change any
operative conditions.

\begin{remark}[Agent neutrality and coupling]\label{rem:coupling}
The definition applies to an individual, business, state, institution, or
coalition. At collective scale, aperture may be distributed across offices and
information systems; Narrative may be carried by institutional memory,
procedures, and commitments; admissibility may be carried by authority and
rules; and meta-analysis may require coordination. These are differences in
representation and cost, not in decision kind.

A collective's partial frames differ by degree and structure of coupling, not
by ontology. Shared representation, communication, authority, binding
procedure, memory, commitment, and coordinated capacity can make the partial
frames more mutually constraining. A so-called group mind is the high-coupling
end of that continuum, not a distinct kind of agent. Where coupling is weak,
the proper object of analysis is an interaction among partial frames; where it
is strong, an effective collective frame can be used at the operative scale.
\end{remark}

\section{Biological evolution: changing the organisation that continues}\label{sec:evolution}

Biological evolution gives a direct example of operative restructuring
without prospective deliberation. Heritable variation changes the structures
organisms carry; differential reproduction changes which structures are
represented in later populations. The resulting organisms realise
continuations through those structures. Their response organisation can
embody the consequences of earlier selection without representing that
selection or reasoning about the future it makes reachable.

At an organism's action scale, inherited and developed organisation conditions
responses. At a population scale, evolutionary change alters the distribution
of such organisation across generations. These are different boundaries and
causal resolutions. Calling an organism's action a realised decision does
not require treating a population as a deliberating agent, or calling each
mutation a decision.

For a supplied sufficient organism description, write
\[
X_{t+1}=T_{\theta}(X_t,B_t),
\]
where \(\theta\) denotes inherited components of response organisation,
\(X_t\) includes relevant developed and current state, and \(B_t\) states
containing conditions. An inherited change \(\theta\to\theta'\) can change
the successor relation, comparison field, physiological feasibility, or
responses under the same environmental cue. Evolutionary change therefore
provides a route to changed operative future geometry without an organism
first selecting a different action in a hypothetical model. The organism's
own lifetime learning and hypothetical processing provide additional routes
to changing determining conditions.

\subsection{Differential continuation without a future model}

Consider a constructed selection-only population with two faithfully inherited
types and a fixed environment. Define \(p_i\) as the initial type frequencies
and \(w_i\) as independently specified reproductive output per type over one
generation. With positive mean output and no mutation, migration, or finite
sampling, the next frequencies are
\[
p'_i=\frac{p_iw_i}{\sum_j p_jw_j}.
\]
For \(p_1=p_2=1/2\) and \(w_1=1,w_2=2\), they are \(1/3,2/3\).
The output definition yields the update by counting descendants and
normalising; no forward comparison by either type is needed. This is a
population bookkeeping instance, not an equation for an organism choosing
its next action. Mutation, variable environments, transmission effects,
and sampling require additional terms in a fuller model. Existing
information-theoretic treatments of selection analyse these distributional
changes \citep{frank2012}; the connection is prior work rather than a new
identification of fitness with informational or physical energy.

\subsection{A measured change in reachable biological futures}

In the long-term \emph{Escherichia coli} experiment, citrate was present in the
medium before a lineage evolved aerobic citrate use. Replay experiments
showed that later sampled genotypes had greater propensity to evolve the
trait than the ancestral and early samples, supporting historical
potentiation \citep{blount2008}. Subsequent genomic analysis identified a
tandem duplication that placed a citrate transporter under an aerobically
expressed promoter \citep{blount2012}. The changed regulatory organisation
enabled a physiological continuation that the earlier organisation did not
implement under those conditions.

In the present framework, the resource was already part of the containing
environment; changed organism structure altered its usable future geometry.
The bacteria did not need to represent citrate use as a possible future and
then approve it. Prior history matters through the structure it leaves and
the subsequent changes that structure permits. The experimental replay and
genomic reconstruction make this more specific than merely attributing an
unobserved preference to the successful lineage.

Evolution is therefore the central example of changed organisation changing
continuation. It also shows why hypothetical processing is one route to
revision rather than its defining requirement. This is distinct from the
momentum paper's maternal-inheritance grounding
(Section~\ref{NM-sec:evolutionary}), whose particular persistence claims
retain their stated limits. Differential reproductive success alone does
not prove that each realised action takes a strict minimum of physical
energy: physiological expenditure, reproductive output, and the operative
action comparison are separately grounded quantities. The minimum rule
below adds that further, independently answerable hypothesis.

\section{Continuation and cost geometry}\label{sec:cost-geometry}

\emph{Seldon}'s Definition~\ref{SM-def:contours} distinguishes local edge cost
\(c_{\Frame}(x,y)\) from accumulated cost-to-reach \(C_{\Frame}(y)\), and
supplies basins only after a successor policy is declared. Neither quantity
is automatically a physical energy, a state potential, or Narrativium.

\begin{definition}[Finite comparison specification]\label{def:finite-policy}
Supply a finite state space, an admissible successor set
\(\operatorname{Succ}_{\Frame}(x)\), and a feasibility predicate
\(\operatorname{Feasible}_{\Frame}(x,y)\) using declared capacity conversions.
Let \(K_{\Frame}(x)\) be the feasible successor set. Declare an effective-cost
field \(q_{\Frame}(x,y)\) with fixed units or comparison convention, and a
deterministic resolution rule selecting one member of each nonempty tied set.
Where \(K_{\Frame}(x)\ne\varnothing\) and every candidate has defined finite
cost, the minimum continuation is
\[
b_{\Frame}(x)=\operatorname{Resolve}_{\Frame}\!\left(
\arg\min_{y\in K_{\Frame}(x)}q_{\Frame}(x,y)\right).
\]
An empty feasible set terminates this policy. A fixed point requires an
explicitly supplied self-loop. If a comparison cost is undefined, this
minimum rule is undefined; no missing value silently becomes zero or infinity.
\end{definition}

One possible grounding is a state score plus a priced transition:
\[
q_{\Frame}(x,y)=U_{\Frame}(y)+\lambda_{\Frame}c_{\Frame}(x,y),
\]
where \(U\), \(c\), and the conversion \(\lambda\) are independently supplied
and commensurable. This formula is a model specification, not an identity
between preference and physical energy. Setting \(U=0\), \(\lambda=1\) gives
the cheapest-successor instance of \emph{Seldon}. A multi-dimensional account
can instead supply a priority order; a scalar field then requires an explicit
conversion. Physical expenditure and comparison score should be separately
reported even when a model relates them.

\begin{hypothesis}[Minimum continuation]\label{hyp:minimum}
Within a declared grounding, action-bearing continuation realises the minimum
of the independently specified effective-cost field over its operative
comparison domain, with the declared tie rule. The rule is instantiated in
the system's dynamics; it need not be recomputed as conscious inference at
each action.
\end{hypothesis}

This gives the proposed ``lowest-energy reachable state in the frame''
statement a precise instance: the next realised state is the lowest defined
effective-cost feasible successor. A history-level comparison over a longer
horizon is a different instance and needs a declared history score and the
mechanism carrying that comparison into action. No claim that local minima
yield the globally cheapest destination follows. If \(q\) is not physical
energy, a joule-denominated interpretation needs additional grounding.

The identity claim places the decision within continuation; the hypothesis
restricts which continuation occurs. Assigning the observed successor cost
zero and all others cost one after observation makes any deterministic
trajectory cheapest and supplies no evidence for the hypothesis. A probability
model permitting nonminimum actions also does not establish an exact minimum
law; a proposed noisy or stochastic extension must state its different rule.

\begin{proposition}[Local minimum need not minimise history cost]\label{prop:local-global}
A minimum-edge successor policy need not realise a minimum-cost history over
a horizon of two transitions.
\end{proposition}
\begin{proof}
Supply exactly two histories \(x\to a\to z\) and \(x\to b\to w\), with
edge costs \(1,10\) and \(2,0\), respectively, and sufficient capacity for
both. Local minimisation chooses \(a\) at \(x\); its history costs \(11\).
The other history costs \(2\). No ties or unknown costs occur.
\end{proof}

Likewise, a cycle under a cheapest-successor policy cannot be strict descent
of one state potential: the potential would have to be strictly smaller upon
returning to the same state. Basins and escape ridges describe the supplied
policy, not a universal gradient of accumulated expenditure. Unresolved
dimensions stay \textsc{null}. Feasibility, membership, and cost remain typed
separately.

The momentum paper's cost-form hypothesis
(Hypothesis~\ref{NM-hyp:narrative-momentum}) supplies positive alteration
cost with independently measured baseline expenditure. Its cheapest-reuse
comparison requires equal baselines and continuing admissibility. A cheaper
baseline can outweigh the alteration charge; a grammar change can remove a
transition altogether. The present account inherits these conditions.

\section{Prospective deliberation and autonomous processing}\label{sec:prospective}

\begin{definition}[Prospective deliberation]\label{def:prospective}
Deliberation is present processing of possible continuations not yet realised
at the declared causal resolution. It may compare, simulate, inhibit,
prepare, or commit, and can change operative conditions of later continuation.
The processing itself is a realised transition with its own costs.
\end{definition}

One can deliberate while acting, or form a commitment about the next causal
step. Reviewing a realised past concerns its present representation and future
consequences. It does not change the realised history. No requirement that
every action be preceded by explicit future modelling follows.

Deliberate cognitive deformation is a particular route to changed conditions.
Where it influences an action or commitment, it supplies the causal basis for
the narrower decision attribution of Definition~\ref{def:decision-attribution}.
Evolutionary restructuring in Section~\ref{sec:evolution} shows that changed
continuation can also arise without that route. A change in behaviour is
therefore insufficient by itself to infer a deliberative decision.

In the familiar System 1/System 2 language, \emph{System 1 simply is}: the
installed autonomous response organisation operates. There is no additional
act of conscious approval required to turn that organisation into action.
System 2 is embedded hypothetical processing that can reorganise conditions
of subsequent action. Both belong to the present operative system.

The research connection is to autonomous defaults and working-memory-dependent
hypothetical processing in \citet{evansstanovich2013}. Their Type 1/Type 2
terminology avoids treating these as exactly two cognitive or neural systems.
This is an interpretive mapping: Type 1 does not name every operative
transition, and Type 2 does not name every revision. Autonomous adaptation
can alter standing structure; hypothetical processing can fail to redirect
it. Processing can overlap, and Type 2 need not always yield a better action.
The cognitive distinction does not prove Hypothesis~\ref{hyp:minimum}.

\subsection{The woodlouse: choice carried in locomotor state}\label{sec:woodlouse}

The woodlouse provides an organism-scale example of physical decision.
In turn-alternation experiments with \emph{Porcellio scaber}, preceding forced
turns influenced subsequent free-turn direction and angle. Hughes's
experiments supported an explanation in terms of bilaterally asymmetrical
leg movements (BALM), with accumulated effects of preceding turns affecting
the subsequent response \citep{hughes1985}. The determining organisation
therefore includes locomotor state carried into the apparent choice point,
rather than just the visible branches of the runway.

In this framework, the next turn is the realised action-bearing continuation
of the woodlouse's body in its containing environment. The preceding movement
changes physical conditions of that continuation. An observer can describe
left and right as alternatives, but that description does not require the
animal to hold an explicit menu of alternatives, understand the route, or
deliberately select a corrective turn. History can constrain action through
the state it leaves in the organism. This is decision in the broad structural
sense of Definition~\ref{def:realised-decision}; it does not by itself warrant
the narrower deliberative attribution.

The apparent choice point makes the provenance issue concrete. A turn by
itself does not disclose whether route comparison or carried locomotor
asymmetry produced it. Controlled changes to preceding movement provide
evidence about the mechanism. The bodily trace is real; identifying its
origin and explanatory role is a model's inference. This is the same
distinction as Remark~\ref{rem:decision-provenance}, applied to physical action.

The body's contribution is part of operative future geometry. Similar
visible junctions need not supply identical operative conditions when
preceding locomotion has left different bodily states. Changing those
conditions changes the subsequent response without requiring a further
chooser outside the response organisation. At this scale, the biological
example connects evolved organisation to physical action, while the
language-model analogy below connects learned organisation to generated
output.

Hughes left unresolved whether the asymmetry's control had a neural basis
or another source such as muscle fatigue. Neural control is itself physical;
the example does not depend on excluding it. The study supports locomotor
asymmetry as a mechanism and requires no demonstrated explicit route
comparison, but does not prove absence of every internal representation or
measure a literal minimum-energy rule. Behavioural competence, mechanism,
explicit understanding, and operative cost remain distinct questions.

\subsection{Competent response does not require explicit understanding}\label{sec:understanding}

The structural implication is that understanding is not an inherent
prerequisite of autonomous action. An installed response can work without
the agent holding an explicit account of why it works, what history installed
it, or how its consequences arise. Evolutionary adaptation, training, and
practice can leave operative structure that produces a response without
retaining a recoverable explanation of that structure's origin. Competence
and explicit causal understanding are therefore separately assessed.

Here explicit causal understanding means a usable model of the relevant
relations and consequences that can support explanation, counterfactual
comparison, and deliberate revision. It is neither guaranteed by successful
output nor excluded by autonomous processing. If understanding instead names
an implicit representation or task-sensitive competence, an autonomous
system can possess it. The proposal is a non-requirement of explicit
understanding, not a definition that strips every Type 1 process of all
semantic or causal representation. Type 2 processing also need not yield
accurate understanding: its represented relations can be incomplete or wrong.

A basic autoregressive language model provides an operational analogy.
Its learned parameters and current context determine a next-token
distribution, and a declared decoding operation produces continuation.
Language-model training can yield task performance from this objective
\citep{radford2019}; generation need not be preceded by a separate explicit
account of why the response is appropriate. Training changes the learned
organisation, while inference uses that organisation with changing context
and activations. The analogy concerns response generation by installed
structure, rather than identity of human and Transformer mechanisms.

The analogy does not establish universal absence of understanding in language
models. A sequence model trained on Othello moves developed a board-state
representation that could be intervened on to change its output
\citep{li2023}. Learned continuation can therefore contain causally effective
internal models even without an explicitly supplied world model. Nor does
sampling from a language model establish the exact minimum policy of
Definition~\ref{def:finite-policy}.

This sharpens the decision-attribution problem. Fluent explanation, competent
action, or a stated reason alone does not establish that an accurate causal
model deliberately deformed operative geometry before the action. Such a
model's existence, accuracy, and influence are different questions. A later
explanation can itself be a generated continuation whose relation to the
earlier action must be investigated. The account permits action without
explicit understanding while leaving particular claims of understanding
answerable to their declared criteria.

Initiating deliberation is itself an operative transition. A scheduled review,
a learned alarm response, or a declared residual threshold can be part of
an instance's installed rules. The account requires no further deliberator
to choose whether to install that rule at each occurrence. Neither the
minimum policy nor the cognitive mapping alone derives a universal trigger.

\section{Narrative continuation and represented decision opening}\label{sec:opening}

The momentum paper defines free continuation as what persists where no rule
writes against standing history-bearing structure (Definition~\ref{NM-def:free-continuation}). At the decision level, this yields a narrower reading.

\begin{definition}[Narrative default]\label{def:narrative-default}
If the continuation induced by \(S_{\Frame}\) remains represented and
admissible, retaining it requires zero \emph{additional decision inference}.
This assigns zero incremental inference charge to reuse of the installed
decision in the declared grammar. Execution, maintenance, and monitoring
can consume resources; delayed or failed implementation is possible.
\end{definition}

\begin{definition}[Decision opening]\label{def:opening}
Within the declared represented grammar, a standing commitment is reopened
only by either a represented surprise
\(\sigma_{\Frame}\), or a current-frame meta-analysis \(m_{\Frame}\):
\[
\operatorname{Change}(\Frame)\ \Longrightarrow\
\Surprise_{\Frame}(\sigma_{\Frame})\ \lor\ \Meta_{\Frame}(m_{\Frame}).
\]
Meta-analysis is an additional decision-inference operation. Its cost and the
conditions that trigger it vary by frame. It is not implementation cost.
\end{definition}
Here the disjunction classifies reopening operations; it does not predict
which meta-analysis trigger occurs. A predictive application must specify,
for example, a scheduled review or a threshold on accumulated residual.
It is not an exhaustive taxonomy of every cause of changed action. External
intervention, capacity loss, or automatic adaptation can change operative
continuation without represented reopening. A surprise can also activate an
installed autonomous response; it is not automatically Type 2 processing.



\begin{proposition}[Direct surprise revision is aperture-resolved]\label{prop:surprise-resolution}
A surprise directly licenses revision only if it has a resolution \(x\) already
available in the current aperture and represented landscape, and that resolution
is admissible and coherent in the current frame:
\[
\operatorname{SurpriseRevision}_{\Frame}(\sigma_{\Frame},x)
\ \Longrightarrow\
x\in L_{\Frame}\text{ within }A_{\Frame}
\ \land\ \Gamma_{\Frame}(x)
\ \land\ \Cohere_{\Frame}(\sigma_{\Frame},x).
\]
\end{proposition}

\begin{proof}
The first conjunct is required by the decision domain
(Definition~\ref{def:domain}); an unrepresented resolution is not yet a
candidate. The second is required by admissibility. The third states that a
resolution cannot borrow its warrant from a successor frame not yet installed.
Without all three, the mismatch can be represented as surprise but not as a
directly selectable replacement. \end{proof}

\subsection{Expected non-occurrence is a negative observation}\label{sec:nonoccurrence}

\begin{definition}[Expected non-occurrence]\label{def:nonoccurrence}
Let \(X\) be a represented event and \(Y\) a represented opportunity condition
or recurrence interval. The frame holds a negative observation \(\neg X\) over
\(Y\) only if all of the following hold:
\begin{enumerate}[label=(\roman*)]
\item \(\Frame\) represents an expectation that \(X\) occurs under \(Y\), such
      as a recurrent relation ``\(X\) each \(Y\)'';
\item \(Y\) is represented as having occurred and the relevant observation
      interval has elapsed;
\item the current aperture can distinguish non-occurrence from an undisclosed
      or unobserved occurrence; and
\item \(X\) is not represented in that interval.
\end{enumerate}
Only then may \(\neg X\) be treated as candidate surprise rather than as
silence or missing information.
\end{definition}

\begin{proposition}[Non-occurrence constrains, but does not identify, change]\label{prop:nonoccurrence}
If an expected non-occurrence satisfies Definition~\ref{def:nonoccurrence}, it
constrains the represented Seldon manifold by making continuations that require
\(X\) under the observed \(Y\) non-coherent with the current observation. It
does not identify a unique hidden change. A probabilistic expectation that
allows absence is instead updated under its declared evidence model; absence
alone does not logically eliminate a history that permits it.
\end{proposition}

\begin{proof}
The represented recurrence and observed opportunity make continuations requiring
\(X\) incompatible with the represented absence. Several distinct represented
causes can remain compatible with that absence: a changed boundary condition, a
changed intervention, a changed coupling, or a changed recurrence. If reliable
non-occurrence cannot be distinguished from observation failure, the
definition's aperture condition is unmet and a weaker evidence model is
required. Hence the absence narrows the
represented future geometry without determining one causal explanation.
\end{proof}

\begin{remark}[Counting what did not happen]\label{rem:nonoccurrence}
A repeated event stopping is therefore not a mere lack of data. Given the
represented expectation and adequate aperture, it is data of negative form. It
may be surprise and thereby open decision. Direct revision still requires the
aperture-available, frame-coherent resolution of
Proposition~\ref{prop:surprise-resolution}; otherwise the absence may motivate
meta-analysis to distinguish the surviving manifold shapes.
\end{remark}

\begin{remark}[Unresolved surprise]\label{rem:unresolved}
An unresolved or frame-incoherent surprise may motivate meta-analysis. That
operation may refine aperture, manifold resolution, active criteria, or the
relation between them. It does not entail that the Narrative changes; it only
pays to reopen comparison.
\end{remark}

\section{Target shapes and feasible selection}\label{sec:targets}

\begin{definition}[Target shape]\label{def:target-shape}
Where \(\Frame\) represents a Seldon manifold, a target is an acceptable shape
\[
\Target_{\Frame,h,\rho}\subseteq\Seldon_{h,\rho}(\Frame),
\]
which may contain regions of future states, families of histories, relations,
or constraints. It is a represented object in \(\Frame\), not a hidden
privileged destination.
\end{definition}
This is the acceptable region of the Seldon theory
(Definition~\ref{SM-def:acceptable}) at the decision frame's horizon.
$T$ is local notation for that region, not an additional kind of object.



\begin{proposition}[A target shape is not a point with an error bar]\label{prop:shape-not-point}
A point estimate \(\hat z\pm\epsilon\) represents uncertainty about one intended
location or value. A target shape represents acceptability: distinct
configurations may be acceptable while nearby configurations may not be. The
two remain different relations even when their represented regions have
the same geometry.
\end{proposition}

\begin{proof}
The first object is indexed by uncertainty around \(\hat z\). The second is
membership in \(\Target\). A disconnected, non-convex, or constraint-defined
\(\Target\) cannot in general be recovered from a point and one scalar radius.
\end{proof}

\begin{definition}[Eligible decision set]\label{def:eligible}
Let
\[
\Eligible_{\Frame}=\left\{x\in L_{\Frame}\ \middle|\
\Gamma_{\Frame}(x)\ \land\ \text{any applicable Narrativium feasibility gate holds}\right\}.
\]
A transition may be desirable yet absent from \(\Eligible_{\Frame}\) because
its represented consequential deformation is not currently reachable.
\end{definition}

\begin{proposition}[Pareto comparison does not select]\label{prop:pareto}
Let \(v_{\Frame}:\Eligible_{\Frame}\to\mathbb{R}^k\) be declared comparison
dimensions and let \(\Pareto_{\Frame}\) be its non-dominated subset. Membership
in \(\Pareto_{\Frame}\) does not select a unique candidate.
\end{proposition}

\begin{proof}
Two candidates can be mutually non-dominated. Then Pareto comparison supplies
no ordering between them. A priority order, scalarisation, lexicographic rule,
or another stated selection relation is still required.
\end{proof}

After a decision opening, selection can ask whether an eligible transition
reaches or remains within \(\Target_{\Frame,h,\rho}\), then apply its stated
selection relation among qualifying histories. Narrativium is still a capacity
gate, not a value of one target point over another. Within this retention
grammar, absent surprise and meta-analysis and with operative conditions
unchanged, the standing commitment remains installed despite visible
non-dominated alternatives. Its persistence is a declared rule; it does not
establish that the installed action minimises any independently supplied
field. A minimum-continuation grounding must separately relate that field
to the installed rule.

\section{Self-decision and operative revision}\label{sec:self}

\begin{proposition}[Successor landscapes are present candidates]\label{prop:self}
A currently represented transformation \(\tau\) may have a successor decision
frame
\[
\Frame' = \tau(\Frame),
\]
with changed aperture, target shape, Narrative, admissibility, or future
manifold. The current frame evaluates \(\tau\) only through what it represents
and orders now.
\end{proposition}

\begin{proof}
Admissibility conditions are functions of antecedent state by the
antecedent-state principle (Proposition~\ref{GE-prop:antecedent-state}). The
newly installed frame cannot authorise the transition that first installs it
(Corollary~\ref{GE-cor:no-self-auth}). Therefore a successor landscape is a
current candidate, not a prior authority.
\end{proof}

\begin{remark}[No homunculus and no future veto]\label{rem:self}
This is a structural constraint, not a claim that every human or artificial
system actually makes self-revisions in this way. It permits direct self-change
under a present frame's rules and capacity, but prevents a successor-only value,
identity, or order from supplying the initial warrant.
\end{remark}

\begin{definition}[Operative revision]\label{def:revision}
A revision changes causally relevant conditions of subsequent continuation:
aperture, horizon, comparison fields, target constraints, standing structure,
admissibility, capacity, coupling, or the transition machinery. A represented
revision is a change in the model held by the agent. It is an operative
revision of action only insofar as it changes conditions generating that
action at the declared scale.
\end{definition}

Thinking of an alternative is a real present process. It need not install
the relations required to realise the alternative. Training, reappraisal,
inhibition, changed incentives, and commitment devices are candidate ways
of changing operative conditions; their efficacy and costs need grounding.
Revising an acceptable region, changing its weight, and changing the set of
physically available actions are different transformations.

\begin{proposition}[Changed continuation requires changed determining conditions]\label{prop:revision-conditions}
In a sufficient deterministic instance, two occurrences with identical
operative state, boundary conditions, and transition rule have identical
successors. A different successor therefore requires a difference in at
least one of those determining conditions.
\end{proposition}
\begin{proof}
The supplied deterministic function \(T\) has one value for fixed arguments.
\end{proof}

This formalises the structural constraint on changing a decision by changing
its operative manifold. It is not independent evidence for minimisation.
Identical visible action menus are weaker than identical determining
conditions: different weights, alarm states, or carried commitments can
produce different actions over the same menu. Conversely, an unexplained
action difference must not be rescued by inventing an unrestricted hidden
field after the event. It challenges the declared grounding or its state
sufficiency.

Revision costs belong to the transition that revises the organisation. A
later cheap continuation does not erase those costs. The present organisation
must carry the trigger, resources, and admissibility of the first revision;
the future self supplies no prior authorisation. This condition does not
derive a uniquely correct response to transformative preference.

\section{Worked finite instances}\label{sec:examples}

The following are constructed instances, not measurements of human energy or
ethical value. They supply domains and costs before applying the rule. All
tables use scalar effective-cost units; physical expenditure is separately
declared where needed. Ties are resolved by the fixed label order \(a<b\).
Action endpoints are terminal unless a further edge is supplied.

\subsection{Installed continuation, expanded aperture, and reappraisal}

At action state \(x\), let the admissible menu be \(\{a,b\}\), with both
actions feasible. The installed response implements
Definition~\ref{def:finite-policy}. Declare the score
\[
q(a)=1+7z,\qquad q(b)=4,
\]
where \(z\in\{0,1\}\) states whether a particular consequence is represented
and coupled into the operative comparison. Both actions have physical
execution expenditure of one resource unit. A review operation costs two
resource units, and the initial resource budget is four. An installed alarm
rule sends the system to review when a cue is registered, before the action
transition. Review sets \(z=1\) and leaves two resource units; without a cue,
the response uses \(z=0\). Thus neither initiating review nor redirecting
action needs an unspecified chooser.

\begin{center}
\begin{tabular}{lrrl}
\toprule
Operative comparison & \(q(a)\) & \(q(b)\) & Next action\\
\midrule
Consequence omitted: \(z=0\) & 1 & 4 & \(a\)\\
Consequence coupled: \(z=1\) & 8 & 4 & \(b\)\\
\bottomrule
\end{tabular}
\end{center}

Without the alarm, the installed response realises \(a\) with zero additional
decision inference and positive execution expenditure. The alarm/review path
realises \(b\) after paying the review charge. The reachable action menu stays
the same; the field changes. The seven score units are not seven physical
resource units. The review rule is supplied as part of this instance, not
derived from the action scores.

This can model expanded aperture when a previously omitted dimension becomes
represented and causally coupled. A separate reappraisal instance keeps that
dimension represented throughout and changes its coefficient from zero to
seven. It produces the same table by changing weighting rather than aperture.
Both descriptions must identify the actual coupling that changes the response.
Neither eliminates consequences that occurred while they were omitted.

\subsection{A represented alternative without operative redirection}

Extend the state with a displayed intention \(p\). The observer knows the
predeclared response rule uses \(z\), not \(p\). A display operation sets
\(p=b\), costs one resource unit, and leaves \(z=0\). The next action remains
\(a\). The display operation changed actual state and represented intention,
but not the field or rule governing the action. This instance distinguishes
knowing or describing an alternative from installing its operative conditions.
It does not assume that intentions never have causal force.

\subsection{Inadmissibility versus limited capacity}

At a fresh \(x\), declare \(q(a)=1,q(b)=4\), and physical action expenditures
\(e(a)=3,e(b)=1\). Set feasibility to \(e(y)\le r\), where \(r\) is current
resource budget. With \(r=3\), both are feasible and the policy realises \(a\).
With \(r=1\), the same admissible menu remains, but only \(b\) is feasible.
Removing the edge to \(a\) by a separate rule change instead makes \(a\)
inadmissible even with \(r=3\). The two deformations can yield the same action
while changing different parts of operative geometry. The budget here is an
explicit physical resource constraint; equating it with a Narrativium gate
would require a separately grounded conversion.

\subsection{Horizon and basins}

The two-transition histories in Proposition~\ref{prop:local-global} have
totals \(11\) and \(2\), although their first edges cost \(1\) and \(2\).
A supplied two-step planning operation can compute both totals and install
the history through \(b\). Give that operation a separate expenditure of one
unit in the same accounting currency, with additive path and planning
expenditure declared for this instance; its planned course then costs \(3\).
Both endpoints are acceptable and both courses are feasible under a budget
of twelve. This construction shows
what a longer-horizon rule can implement, but does not derive its initiation
or price from the local rule.

For a basin example, separately supply edges \(u\to v\), \(v\to u\), and
\(v\to w\) with local costs \(1,1,3\), and a declared absorbing self-loop at
\(w\) of cost zero. Under the local policy, \(\{u,v\}\) is an invariant cycle
with basin \(\{u,v\}\). Its minimum escape-edge cost is three, but that edge
is never taken by the unchanged cheapest-successor policy. An external
intervention, constraint revision, or changed comparison field is needed to
produce escape. The cycle is an admissible instance, not descent to a global
minimum of a single state potential.

\section{Decision quality, implementation, and outcome}\label{sec:quality}

The account is not an ideal-decision theory. In the retention grammar, an
unchanged operative frame carries its installed commitment without reopening.
A frame-valid but poor represented selection can cohere within a limited or
miscalibrated aperture. An agent may pay for meta-analysis and still fail to
realise the intended action. Implementation may fail while the represented
selection or commitment remains unchanged. An intended action may also be
realised yet lead to undesirable consequences. These are distinct diagnosis
classes; an unchanged commitment does not make a failed intended action into
the same realised decision.

This distinction keeps outcome learning correctly typed. A later consequence may
become represented surprise, which can then have an aperture-available,
frame-coherent resolution or motivate meta-analysis. It does not make the
external consequence a retroactive input to the original decision.

Whole-system consequences are also distinct from the operative comparison
that generated action. The informational accounting companion
(Section~\ref{RI-sec:decisions}) separates an actual consequence account from
the model and relative weights through which an agent evaluates it.
An omitted casualty, future loss, or burden on another system can occur
regardless of the actor's representation. Changes to the actor's own
organisation, including moral injury and subsequent self-modelling, also
belong to the consequence account. They can alter later operative fields
without making the original internal score identical to whole-system cost.
Action can be a minimum under a grounded operative field and still cause
large externally measurable losses.

\section{Operational commitments and research neighbours}\label{sec:operational}

The operational commitment is comparative. For a stated population, frame
class, target shape, horizon, baseline, and inference budget, the grammar must
be tested for:
\begin{enumerate}[label=(\arabic*)]
\item predictive value for persistence, surprise, meta-analysis, revision, and
      decision-relevant consequences;
\item decision value for target-shape membership or other declared objectives;
\item cost of representation, review, and meta-analysis; and
\item transfer limits across aperture, coupling, and horizon.
\end{enumerate}
A failure to improve these relative to a baseline limits operational
usefulness in that scope. A repeated nonminimum response under independently
grounded sufficient state, fields, constraints, and rules challenges the
minimum-continuation hypothesis for that grounding. The identity claim alone
does not distinguish such responses and must not substitute for this test.

\subsection{A discriminating proposed study}

Hold the available action menu fixed and specify an alarm/review operation
that changes an independently measured consequence representation or its
operative weighting. Fit a declared field, trigger rule, revision expenditure,
and state description on a calibration set, then freeze them. On held-out
cases, predict persistence without effective revision and a switch where
revision reverses the feasible ordering. Compare the prediction with a
fixed-frame model that has no revision term, accounting for the extra
parameters and inference expenditure.

The proposal needs a manipulation check that measures the specified operative
change independently of the target action; a verbal statement of intention
alone is insufficient. If the field crosses the preregistered boundary but
the response repeatedly fails to switch under the declared conditions, that
counts against the grounding. A probabilistic observation model can specify
measurement uncertainty in advance; it must not assign arbitrary fields to
each failure. A demonstration only that attention changes behaviour would
not uniquely support this theory. Exact recovery of a brain mechanism is
not required to propose the principle, but evidence must constrain possible
groundings and distinguish the claimed rule from competing accounts.

\subsection{Research neighbours and contribution}

Dual-process cognition supplies the distinction between autonomous responses
and hypothetical processing discussed in Section~\ref{sec:prospective}
\citep{evansstanovich2013}. The present account interprets both as embedded
continuations and specifies how hypothetical processing can alter operative
future geometry. That connection does not establish a universal brain
architecture or minimum-cost dynamics.

Several other literatures are close but non-equivalent. Resource-rational
metareasoning treats additional reasoning as costly and models selection of
reasoning processes \citep{godara2024}; it does not define Narrative or this
opening condition. Default research establishes heterogeneous effects of
preselected options \citep{jachimowicz2019}; it does not establish a general
history-bearing Narrative law. Active inference supplies a Bayesian policy model
using prediction and preference models \citep{dacosta2024}; GE residual is not
thereby variational free energy, prediction error, or Bayesian surprise.
Multiobjective optimisation confirms that Pareto-optimal alternatives require
additional preference or selection information for final choice
\citep{pajasmaa2025}. AGM belief revision provides a formal neighbour for
retention and change of commitments \citep{agm1985}; it is not identical to
Narrative revision.

The contribution is the integration of decision as embedded continuation,
represented futures, and paid restructuring within the series' framework.
The finite policy and examples expose the supplied premises. They claim no
priority over planning, metareasoning, dynamical accounts of cognition, or
dual-process theory, and derive no general neural or substrate-independent
energy law from those neighbouring results.

\section{Limits}\label{sec:limits}

The paper proposes an ontological interpretation and supplies implementable
finite dynamics under declared premises. It does not establish that a nervous
system, organisation, country, or artificial agent instantiates the complete
grammar or the exact minimum rule. The interpretation requires neither
conscious choice nor a complete self-model; its explanatory value depends on
grounding more restrictive dynamics than the identity alone supplies.

No universal scalar energy, global optimisation theorem, interpersonal
Narrativium accounting, target shape, meta-analysis trigger, or neural mechanism
is derived. The local construction needs known comparison costs, nonempty
feasible domains, and declared tie rules. Partial observations cannot warrant
complete-state predictions. A represented future, another agent's model, and
actual operative organisation remain distinct. At collective scale the account
uses declared coupling among partial frames, rather than a categorical group
mind.

The central proposal is therefore precise but conditional: realised decision
is action-bearing continuation; deliberation is present processing of
unrealised futures; effective redirection changes operative conditions.
Minimum continuation adds an independently answerable hypothesis. Its human
application remains open to bounded comparative investigation.

\begin{thebibliography}{99}

\bibitem[Alchourr\'on et~al.(1985)]{agm1985}
Carlos E. Alchourr\'on, Peter G\"ardenfors, and David Makinson.
On the logic of theory change: Partial meet contraction and revision functions.
\emph{Journal of Symbolic Logic}, 50(2):510--530, 1985.

\bibitem[Blount et~al.(2008)]{blount2008}
Zachary D. Blount, Christina Z. Borland, and Richard E. Lenski.
Historical contingency and the evolution of a key innovation in an
experimental population of \emph{Escherichia coli}.
\emph{Proceedings of the National Academy of Sciences}, 105(23):7899--7906, 2008.
\href{https://doi.org/10.1073/pnas.0803151105}{doi:10.1073/pnas.0803151105}.

\bibitem[Blount et~al.(2012)]{blount2012}
Zachary D. Blount, Jeffrey E. Barrick, Carla J. Davidson, and Richard E. Lenski.
Genomic analysis of a key innovation in an experimental \emph{Escherichia coli}
population.
\emph{Nature}, 489:513--518, 2012.
\href{https://doi.org/10.1038/nature11514}{doi:10.1038/nature11514}.

\bibitem[da Costa et~al.(2024)]{dacosta2024}
Lancelot da Costa, Samuel Tenka, Dominic Zhao, and Noor Sajid.
Active inference as a model of agency.
\emph{arXiv preprint arXiv:2401.12917}, 2024.

\bibitem[Evans and Stanovich(2013)]{evansstanovich2013}
Jonathan St. B. T. Evans and Keith E. Stanovich.
Dual-process theories of higher cognition: Advancing the debate.
\emph{Perspectives on Psychological Science}, 8(3):223--241, 2013.
\href{https://doi.org/10.1177/1745691612460685}{doi:10.1177/1745691612460685}.

\bibitem[Frank(2012)]{frank2012}
Steven A. Frank.
Natural selection. V. How to read the fundamental equations of evolutionary
change in terms of information theory.
\emph{Journal of Evolutionary Biology}, 25(12):2377--2396, 2012.
\href{https://doi.org/10.1111/jeb.12010}{doi:10.1111/jeb.12010}.

\bibitem[Godara et~al.(2024)]{godara2024}
Prakhar Godara, Tilman Diego Alem\'an, and Angela J. Yu.
Metareasoning in uncertain environments: A meta-BAMDP framework.
\emph{arXiv preprint arXiv:2408.01253v1}, 2024.
\url{https://arxiv.org/abs/2408.01253v1}.

\bibitem[Hughes(1985)]{hughes1985}
R. N. Hughes.
Mechanisms for turn alternation in woodlice (\emph{Porcellio scaber}):
The role of bilaterally asymmetrical leg movements.
\emph{Animal Learning \& Behavior}, 13(3):253--260, 1985.
\href{https://doi.org/10.3758/BF03200018}{doi:10.3758/BF03200018}.

\bibitem[Jachimowicz et~al.(2019)]{jachimowicz2019}
Jon M. Jachimowicz, Shannon Duncan, Elke U. Weber, and Eric J. Johnson.
When and why defaults influence decisions: A meta-analysis of default effects.
\emph{Behavioural Public Policy}, 3(2):159--186, 2019.

\bibitem[Johnson et~al.(1993)]{johnson1993}
Marcia K. Johnson, Shahin Hashtroudi, and D. Stephen Lindsay.
Source monitoring.
\emph{Psychological Bulletin}, 114(1):3--28, 1993.
\href{https://doi.org/10.1037/0033-2909.114.1.3}{doi:10.1037/0033-2909.114.1.3}.

\bibitem[Li et~al.(2023)]{li2023}
Kenneth Li, Aspen K. Hopkins, David Bau, Fernanda Vi\'egas,
Hanspeter Pfister, and Martin Wattenberg.
Emergent world representations: Exploring a sequence model trained on a
synthetic task.
\emph{International Conference on Learning Representations}, 2023.
\url{https://arxiv.org/abs/2210.13382}.

\bibitem[Pajasmaa et~al.(2025)]{pajasmaa2025}
Juuso Pajasmaa, Kaisa Miettinen, and Johanna Silvennoinen.
Group decision making in multiobjective optimization: A systematic literature review.
\emph{Group Decision and Negotiation}, 34:329--371, 2025.

\bibitem[Radford et~al.(2019)]{radford2019}
Alec Radford, Jeffrey Wu, Rewon Child, David Luan, Dario Amodei,
and Ilya Sutskever.
Language models are unsupervised multitask learners.
\emph{OpenAI technical report}, 2019.
\url{https://cdn.openai.com/better-language-models/language_models_are_unsupervised_multitask_learners.pdf}.

\end{thebibliography}

\end{document}
